\documentclass[10pt,twoside,a5paper]{article}
%\usepackage[latin1]{inputenc}
%\usepackage[utf8]{inputenc}
\usepackage[applemac]{inputenc}
\usepackage[T1]{fontenc} % permet les guillemets
\usepackage{amssymb,amsmath}
\def\volnum{Manuscrit}
%\def\volnum{Volume 49, 2025}
\usepackage{graphicx}
\usepackage{epstopdf} 
\DeclareGraphicsRule{.tif}{png}{.png}{`convert #1 `basename #1 .tif`.png} 
%\usepackage[english,french]{babel}
%\selectlanguage{french}

% Format des AFLB
%\usepackage{aflb250330}
\usepackage{aflb-web}
\pdebut{1}
\def\pacs#1{\LP P.A.C.S.: #1}
\def\email#1{{email: \tt #1}}
\title{Titre complet, plusieurs lignes si n\'ecessaire}
\titleshort{Titre court \dots}
\author{Auteur}
\authorshort{Auteur}
\address{Adresse}

\setlength{\arraycolsep}{1pt}

\begin{document}
\maketitle

\vskip 1cm
\begin{abstract}
R\'ESUM\'E. La r\'edaction des Annales pr\'esente dans cette note l'objet de la
publication, les conditions de soumission des articles, les diff\'erents types
d'articles. Le fichier de style n\'ecessaire, AFLBCOURS.STY est disponible sur le
site de la Fondation Louis de Broglie : 

\hfil
{\tt http:/www.FondationLouisdeBroglie.org},\hfil

 ou par courrier \'electronique \`a : 
{\tt fond.broglie@wanadoo.fr}.

{\it
ABSTRACT. The editors present in this note the object of the publication, the
conditions of contribution of the articles, the various types of articles.
Most of the necessary style file AFLBCOURS.STY is available on our website:
{\tt http:/www.FondationLouisdeBroglie.org} or on request by electronic mail to:
{\tt fond.broglie@wanadoo.fr}.}
\end{abstract}
\pacs{03.65.Bz; 04.20.-q}


\section{L'objet et les limites}

Les Annales de la Fondation Louis de Broglie ont pour objet la publication de
travaux relatifs \`a la physique th\'eorique, et en particulier \`a la m\'ecanique
quantique et ses applications aussi bien th\'eoriques qu'exp\'erimentales. Les Annales
encouragent les articles de revue, les courtes notes et les Lettres \`a l'Editeur. 


Les manuscrits soumis pour publication sont \`a adresser dans la forme de ces
instructions \`a : 

\hfil{Fondations Louis de Broglie 23, rue Marsoulan, F-75012 PARIS}
\hfil
 
\LP ou, de pr\'ef\'erence, par courrier \'electronique \`a :

\hfil{\tt fond.broglie@wanadoo.fr} \hfil

\LP au
format des Annales d\'ecrit dans ces instructions. Les logiciels que les Annales
peuvent exploiter sont \TeX, ScientificWord et Word de Microsoft. {\it Notre
pr\'ef\'erence va \`a \LaTeX{} (ou \TeX{}) et ScientificWord}. En effet la grande
vari\'et\'e des
\'editions du logiciel de Microsoft avec les probl\`emes de compatibilit\'e entre
elles, nous conduit \`a limiter le plus possible l'usage de ce traitement de
texte. Le fichier Word ``{\tt AFLB-Instructions.doc}" contient tous les styles
n\'ecessaires \`a Word, associ\'es \`a la langue choisie.

Avant publication un transfert de copyright de ou des auteurs \`a la Fondation
Louis de Broglie devra \^etre r\'egularis\'e.  

La longueur des articles ne d\'epassera
pas 6 000 mots sauf accord exceptionnel avec la r\'edaction.  Les langues admises
sont le fran\c cais et l'anglais. Le titre et un court r\'esum\'e d'environ 100 mots
devront \^etre soumis dans les deux langues.

\section{Mise en page: valeurs d\'efinies dans le fichier de style}

La page imprim\'ee est d\'efinie par sa hauteur : 17,7cm, sa largueur : 11 cm.

Les instructions de mise en page sont donn\'ees pour un format de papier A4 soit
210mm par 297mm.  Le pav\'e de texte incluant les en-t\^etes commence \`a 1,00 cm du
haut de la page. La marge du haut est donc de 1,00 cm celle du bas de 11,00 cm. 

La marge de gauche est de 2,54 cm, en position portrait celle de droite est de
7,46 cm. Pour des raisons de mise en page la r\'edaction des Annales peut \^etre
amen\'ee \`a modifier tr\`es l\'eg\`erement la hauteur de la page

Le titre de l'article a pour seules Majuscules la premi\`ere lettre et
\'eventuellement celles li\'ees aux usages comme dans les formules de la chimie ex :
Al$_2$O$_3$.  Il commence \`a 2 cm de l'En-t\^ete de la 1\`ere page.


\subsection {R\'ef\'erences} 

Elles sont regroup\'ees \`a la fin de l'article. Nous utilisons une
liste num\'erot\'ee avec retrait n\'egatif de 0,7 cm. Pour un livre voir la r\'ef\'erence
[1], pour un article tel celui de Pierre Curie sur la sym\'etrie 
voir la r\'ef\'erence [2], ou l'un des articles les plus c\'el\`ebres de L.
de Broglie \cite{Broglie2}%
\footnote{On pourra utiliser commme ici la num\'erotation automatique propos\'ee par
LaTex.}
 etc. 

\subsection{Notes de bas de pages} 

Elles doivent \^etre exceptionnelles et courtes. {\it Les notes longues sont
souvent mieux en annexes et seront mises comme telles en fin de texte juste avant
les r\'ef\'erences}.

\subsection{Figures}

Les figures seront de pr\'ef\'erence fournies sous forme de fichiers de type EPS
(encapsulated postscript). Elles peuvent \^etre plac\'ees indiff\'eremment dans le
texte, ou en annexes en fin de l'article, \`a l'aide des macros standard des
``packages" de type GRAPHICX.



\vskip 30pt
\begin{frenchref}
% eref si anglais
\bibitem{} L. de Broglie, {\it La th\'eorie de la mesure en m\'ecanique
ondulatoire}, Paris, Gauthier-Villars, 1957.
\bibitem{}Curie P., J. de Phys., 3-i\`eme s\'erie, {\bf 3}, 393-415, (1894)

\bibitem{Broglie2} L. de Broglie, C. R. Ac. Sc.  {\bf 177}, 517, (1923).
\end{frenchref}

\man{27 mai 2003}
\end{document}
